\chapter{Describing Data}\label{ch:b1:describing-data}

In 1978, 2,490 American men of a household panel earned on average 21,554
dollars of 1982, and $51.4\%$ of them earned less than that average. (The
samples of this chapter and the next three measure earnings in dollars of
1982.) % ledger: b1:nsw-1982dollars
% computed: code/bachelor-1/29-describing-data/outputs/histogram.txt
In the United States in 2022, the richest hundredth of the population had
an average income 9.37 times the median. % figdata: pip_stats
A mean is a fine summary of a symmetric distribution and a poor one of a
skewed one, and incomes are skewed. This chapter describes data before any
model: distributions, their centre and spread, how two variables move
together, how means differ across groups --- and the traps in each.

\begin{openingchart}{Average income per person of each hundredth of the
population, United States, 2022, dollars a day at 2021 purchasing power, on
an ordinary axis (left) and a logarithmic axis (right).}{World Bank, Poverty
and Inequality Platform (\omterm{def:b1:describing-data:quantile}{percentiles})}
\begin{tikzpicture}
\begin{groupplot}[omrecord, group style={group size=2 by 1, horizontal sep=1.6cm}, width=0.46\linewidth,
  height=0.36\linewidth, xmin=0, xmax=101, xlabel={Hundredth of the population}]
  \nextgroupplot[ylabel={Dollars a day}, ymin=0, ymax=650]
  \addplot[omDef, thick, sharp plot] table[x=percentile, y=avg_welfare, col sep=comma] {figdata/bachelor-1/29-describing-data/out/pip_usa.csv};
  \nextgroupplot[ymode=log, ymin=0.2, ymax=1000, log basis y=10]
  \addplot[omDef, thick, sharp plot] table[x=percentile, y=avg_welfare, col sep=comma] {figdata/bachelor-1/29-describing-data/out/pip_usa.csv};
\end{groupplot}
\end{tikzpicture}
\end{openingchart}

On the ordinary axis the bottom four-fifths are squeezed together and the
top hundredth stands alone at 629 dollars a day. On the logarithmic axis,
equal distances are equal ratios: the median, 67 dollars, is 2.7 times the
tenth hundredth and 0.4 times the ninetieth. % figdata: pip_stats

\begin{recall}
Mean, median, quartiles, variance and standard deviation were met in the
mathematics series; the Lorenz curve, the Gini coefficient and top
shares at school. Python and pandas are \cref{ch:b1:economic-data-python}.
\end{recall}

\begin{measure}[Income in a survey]
The opening chart's incomes come from a household survey. Household income
is divided by the number of members, which assumes that everyone in the
household shares equally; and wealthy households are less likely to answer,
so the distribution may have too thin a tail at the top. % ledger: b1:pip-methods
Tax records, which reach the top better, give larger top shares.
\end{measure}

\section{Histograms and quantiles}\label{sec:b1:describing-data:histograms}

\begin{definition}[Histogram]\label{def:b1:describing-data:histogram}
\index{histogram}
A \emph{histogram} divides the range of a variable into bins and draws over
each a bar whose area is the share of observations in the bin: its height
is the share divided by the bin's width, a \emph{density}, so that the areas
sum to 1 whatever the widths.
\end{definition}

\begin{definition}[Empirical distribution function]\label{def:b1:describing-data:edf}
\index{empirical distribution function}
The \emph{empirical distribution function} of a sample $x_1, \dots, x_n$ is
$\hat F(t) = \frac1n\#\{i : x_i \le t\}$, the share of observations at or below $t$.
\end{definition}

\begin{definition}[Quantile, percentile]\label{def:b1:describing-data:quantile}
\index{quantile}\index{percentile}
The \emph{quantile} of order $p$ is the smallest value $t$ with $\hat F(t) \ge p$ (or an
interpolation between two observations, as pandas does by default); the
$p$-th \emph{percentile} is the quantile of order $p/100$.
\end{definition}

\omcode{code/bachelor-1/29-describing-data/listings/histogram.py}{1}{17}{A \omterm{def:b1:describing-data:histogram}{histogram} as a table of
densities, and five \omterm{def:b1:describing-data:quantile}{quantiles}.}\label{lst:b1:describing-data:histogram}

\omoutput{code/bachelor-1/29-describing-data/outputs/histogram.txt}{A long right tail: a few men
above 100,000 dollars, a tenth at zero.}\label{out:b1:describing-data:histogram}

\section{Mean, variance and skewness}\label{sec:b1:describing-data:moments}

\begin{definition}[Skewness]\label{def:b1:describing-data:skewness}
\index{skewness}
The \emph{skewness} of a sample is $\frac1n\sum_i (x_i - \bar x)^3/s^3$, with $s^2 = \frac1n\sum_i (x_i - \bar x)^2$:
positive when the right tail is longer, negative when the left one is,
zero for a symmetric distribution.
\end{definition}

\omcode{code/bachelor-1/29-describing-data/listings/moments.py}{1}{13}{The first three moments, of
earnings and of their logarithm.}\label{lst:b1:describing-data:moments}

\omoutput{code/bachelor-1/29-describing-data/outputs/moments.txt}{Positive earnings skewed to the right
($1.645$); their logarithm skewed to the left ($-1.913$).}\label{out:b1:describing-data:moments}

The logarithm pulls in the long right tail, but it stretches the left one:
a man with very low positive earnings sits far below the others on a log
scale. In Brenland's survey, where no household has income near zero, the
logarithm of \omterm{def:b1:national-accounts-in-depth:gni}{disposable income} has a \omterm{def:b1:describing-data:skewness}{skewness} of $-0.123$, against $1.436$ for
the level (weekend problem). % computed: code/bachelor-1/29-describing-data/exercises/solutions.py
Whether the logarithm makes a distribution symmetric depends on its bottom
as much as on its top.

\section{Covariance and correlation}\label{sec:b1:describing-data:correlation}

\begin{definition}[Sample covariance, correlation coefficient]\label{def:b1:describing-data:correlation}
\index{covariance, sample}\index{correlation coefficient}
The \emph{sample covariance} of $x$ and $y$ is $s_{xy} = \frac1{n-1}\sum_i (x_i - \bar x)(y_i - \bar y)$; the
\emph{correlation coefficient} is $r = s_{xy}/(s_x s_y)$.
\end{definition}

\begin{proposition}[Range of the correlation]\label{prop:b1:describing-data:range}
$-1 \le r \le 1$, with $|r| = 1$ if and only if the points lie on a straight line
that is not horizontal or vertical.
\end{proposition}

\begin{proof}
Let $a_i = x_i - \bar x$ and $b_i = y_i - \bar y$. By the Cauchy--Schwarz inequality,
$(\sum_i a_ib_i)^2 \le \sum_i a_i^2 \sum_i b_i^2$, so $s_{xy}^2 \le s_x^2 s_y^2$ and $r^2 \le 1$. Equality holds when
$b = \lambda a$ for some $\lambda \ne 0$, that is $y_i - \bar y = \lambda(x_i - \bar x)$ for every $i$.
\end{proof}

\omcode{code/bachelor-1/29-describing-data/listings/covariance.py}{1}{17}{Covariance and correlation
by hand and with pandas, and a correlation that misses a parabola.}\label{lst:b1:describing-data:covariance}

\omoutput{code/bachelor-1/29-describing-data/outputs/covariance.txt}{Earnings in 1975 and 1978:
$r = 0.73$. A parabola: $r$ near 0.}\label{out:b1:describing-data:covariance}

A correlation measures how close the points are to a straight line. It
says nothing about a relation that is not linear: on a symmetric range, $v = u^2$
is entirely determined by $u$, and the correlation is near zero.

\begin{omfigure}
\begin{tikzpicture}
\begin{groupplot}[omscatter, group style={group size=3 by 1, horizontal sep=0.9cm}, width=0.34\linewidth,
  height=0.3\linewidth, xticklabels={}, yticklabels={}]
  \nextgroupplot[title={\small $r = 0.995$}]
  \addplot[only marks, mark=*, mark size=0.8pt, omDef] table[x=x1, y=y1, col sep=comma] {figdata/bachelor-1/29-describing-data/out/clouds.csv};
  \nextgroupplot[title={\small $r = 0.62$}]
  \addplot[only marks, mark=*, mark size=0.8pt, omDef] table[x=x2, y=y2, col sep=comma] {figdata/bachelor-1/29-describing-data/out/clouds.csv};
  \nextgroupplot[title={\small $r = 0.04$}]
  \addplot[only marks, mark=*, mark size=0.8pt, omDef] table[x=x3, y=y3, col sep=comma] {figdata/bachelor-1/29-describing-data/out/clouds.csv};
\end{groupplot}
\end{tikzpicture}

{\small Three seeded samples of 200 points: close to a line, a cloud, and a
parabola whose correlation is near zero although $y$ is a function of $x$.
% figdata: clouds_r
}
\end{omfigure}

\section{Conditional means}\label{sec:b1:describing-data:conditional}

\begin{definition}[Conditional mean (in data)]\label{def:b1:describing-data:conditional-mean}
\index{conditional mean}
The \emph{conditional mean} of $y$ given $x = a$ in a sample is the mean of $y$ over
the observations with $x = a$; as a function of $a$, it describes how the average
of $y$ changes with $x$.
\end{definition}

\omcode{code/bachelor-1/29-describing-data/listings/group_means.py}{1}{8}{Mean and median earnings
by years of schooling.}\label{lst:b1:describing-data:groups}

\omoutput{code/bachelor-1/29-describing-data/outputs/group_means.txt}{Earnings rise with
schooling, with a dip at 11 years and a step at 12.}\label{out:b1:describing-data:groups}

\begin{omchart}{Mean earnings in 1978 by years of schooling, PSID comparison
sample: the 2,307 men with 8 to 17 years of schooling (counts above the
bars).}{LaLonde (1986) and Dehejia and
Wahba (1999) samples of the National Supported Work data}
\begin{tikzpicture}
\begin{axis}[omrecord, ybar, bar width=9pt, xlabel={Years of schooling}, ylabel={Mean earnings (dollars of 1982)},
  xmin=7.3, xmax=17.7, ymin=0, ymax=38000, width=0.62\linewidth, height=0.38\linewidth, xtick={8,...,17},
  nodes near coords, point meta=explicit symbolic, nodes near coords style={font=\tiny}]
  \addplot[fill=omDef!60, draw=omDef] table[x=education, y=mean, meta=count, col sep=comma] {figdata/bachelor-1/29-describing-data/out/school.csv};
\end{axis}
\end{tikzpicture}
\end{omchart}

The samples were assembled to test whether non-experimental methods can
reproduce the result of an experiment (LaLonde 1986; Dehejia and Wahba
1999). % ledger: b1:lalonde1986
% ledger: b1:dehejia1999
The \omterm{def:b1:describing-data:conditional-mean}{conditional means} describe; whether a year of schooling causes the
step from $16{,}153$ to $20{,}564$ dollars between 11 and 12 years is the question of
\cref{ch:b1:simple-regression}.
% computed: code/bachelor-1/29-describing-data/outputs/group_means.txt

\section{Traps}\label{sec:b1:describing-data:traps}

\begin{definition}[Composition effect, Simpson's paradox]\label{def:b1:describing-data:simpson}
\index{composition effect}\index{Simpson's paradox}
A \emph{composition effect} is a change or difference in an overall mean that
comes from the shares of the groups rather than from the means within
them. \emph{Simpson's paradox} is its extreme case: a comparison that goes one
way within every group and the other way overall.
\end{definition}

\omcode{code/bachelor-1/29-describing-data/listings/simpson.py}{1}{12}{Older and younger households
in Brenland's survey, overall and within working and non-working households.}\label{lst:b1:describing-data:simpson}

\omoutput{code/bachelor-1/29-describing-data/outputs/simpson.txt}{Older households have more
income within each group, and less overall.}\label{out:b1:describing-data:simpson}

\begin{omfigure}
\begin{tikzpicture}
\begin{axis}[omscatter, xmin=0.5, xmax=3.5, ymin=5000, ymax=35000, xtick={1,2,3},
  xticklabels={Not working, Working, All}, ylabel={Mean disposable income (dalens)}, width=0.56\linewidth,
  height=0.38\linewidth, legend style={font=\scriptsize, at={(0.02,0.98)}, anchor=north west}]
  \addplot[only marks, mark=*, mark size=2pt, omDef] table[x=pos, y=young, col sep=comma] {figdata/bachelor-1/29-describing-data/out/simpson.csv};
  \addplot[only marks, mark=square*, mark size=2pt, omThm] table[x=pos, y=old, col sep=comma] {figdata/bachelor-1/29-describing-data/out/simpson.csv};
  \legend{Head under 50, Head 50 or over}
\end{axis}
\end{tikzpicture}

{\small \omterm{def:b1:describing-data:simpson}{Simpson's paradox} in Brenland's survey: within non-working and within
working households, households with older heads have more income; overall
they have less, because $51.7\%$ of them do not work, against $5.6\%$ of the others.}
\end{omfigure}

\begin{record}[Admissions by department]
In 1973 the six largest departments of a graduate school admitted $44.5\%$ of
2,691 male applicants and $30.4\%$ of 1,835 female applicants. % ledger: b1:bickel1975
Department by department the gap largely disappeared: women had applied
more often to the departments that admitted a small share of all
applicants. The pooled figures showed a pattern that the departments'
own decisions did not (Bickel, Hammel and O'Connell 1975).
\end{record}

\begin{definition}[Top-coding, censoring]\label{def:b1:describing-data:top-coding}
\index{top-coding}\index{censoring}
A variable is \emph{top-coded} when values above a cap are recorded at the cap.
More generally, it is \emph{censored} when values beyond a limit are known only
to lie beyond it.
\end{definition}

\omcode{code/bachelor-1/29-describing-data/listings/top_coding.py}{1}{15}{A quarter of a sample at the
cap.}\label{lst:b1:describing-data:top-coding}

\omoutput{code/bachelor-1/29-describing-data/outputs/top_coding.txt}{4,145 men recorded at 25,564.70
dollars: the recorded mean is a lower bound.}\label{out:b1:describing-data:top-coding}

\begin{omchart}{Earnings in 1978 in dollars of 1982, density per 100,000
dollars, bins of 2,000 dollars: PSID comparison sample (blue) and CPS comparison sample (red),
values above 60,000 put in the last bin.}{National Supported Work comparison
samples}
\begin{tikzpicture}
\begin{axis}[omrecord, xlabel={Earnings in 1978 (dollars of 1982)}, ylabel={Density}, xmin=0, xmax=60000, ymin=0,
  width=0.66\linewidth, height=0.36\linewidth, scaled x ticks=false]
  \addplot[omDef, thick, const plot mark mid] table[x=lo, y=psid, col sep=comma] {figdata/bachelor-1/29-describing-data/out/earn_hist.csv};
  \addplot[omThm, thick, const plot mark mid] table[x=lo, y=cps, col sep=comma] {figdata/bachelor-1/29-describing-data/out/earn_hist.csv};
\end{axis}
\end{tikzpicture}
\end{omchart}

The CPS sample's spike at its cap holds a quarter of the men. Every mean,
variance or regression computed on it treats 25,564.70 dollars as their
earnings.

\begin{definition}[Outlier, trimmed mean]\label{def:b1:describing-data:outlier}
\index{outlier}\index{trimmed mean}
An \emph{outlier} is an observation far from the bulk of the data. The
\emph{$\alpha$-trimmed mean} drops the fraction $\alpha$ of the observations at each end and
averages the rest.
\end{definition}

\omcode{code/bachelor-1/29-describing-data/listings/trimmed.py}{1}{13}{The mean moves with one
observation; the median and the \omterm{def:b1:describing-data:outlier}{trimmed mean} barely.}\label{lst:b1:describing-data:trimmed}

\omoutput{code/bachelor-1/29-describing-data/outputs/trimmed.txt}{Removing one man in 445 lowers the
mean by 124 dollars and the median by 29.}\label{out:b1:describing-data:trimmed}

\section*{Exercises}\label{sec:b1:describing-data:exercises}

\begin{exercise}[$\star$]\label{exo:b1:describing-data:read}
From the opening chart's table, the median hundredth of Americans had 67.15
dollars a day in 2022 and the ninetieth 169.65. % figdata: pip_stats
Give the ratio, and say where each lies on the logarithmic panel.
\end{exercise}

\begin{exercise}[$\star$]\label{exo:b1:describing-data:five}
Find the mean and the median of $2$, $3$, $3$, $4$, $18$. Which describes the typical value?
\end{exercise}

\begin{exercise}[$\star$]\label{exo:b1:describing-data:r}
Two variables have a covariance of 12 and standard deviations of 3 and 5.
Find their correlation.
\end{exercise}

\begin{exercise}[$\star$]\label{exo:b1:describing-data:trap}
A country's average wage falls in a year in which the wage of every
occupation rises. Name the trap and give a mechanism.
\end{exercise}

\begin{exercise}[$\star\star$]\label{exo:b1:describing-data:log}
For the PSID men earning at least 5,000 dollars, compute the \omterm{def:b1:describing-data:skewness}{skewness} of
earnings and of their logarithm, and compare with \cref{out:b1:describing-data:moments}.
\end{exercise}

\begin{exercise}[$\star\star$]\label{exo:b1:describing-data:cap}
Replace the CPS sample's capped values by 1.5 times the cap. By how much does
the mean change? And the median?
\end{exercise}

\begin{exercise}[$\star\star$]\label{exo:b1:describing-data:step}
From \cref{out:b1:describing-data:groups}, give the step in mean earnings from 11 to 12
years of schooling, and the step from 12 to 16 years per year.
\end{exercise}

\begin{exercise}[$\star\star$]\label{exo:b1:describing-data:prove}
Prove \cref{prop:b1:describing-data:range} without the Cauchy--Schwarz inequality, using
$\sum_i (a_i/\sqrt{\sum a^2} \pm b_i/\sqrt{\sum b^2})^2 \ge 0$.
\end{exercise}

\begin{exercise}[$\star\star$]\label{exo:b1:describing-data:billionaire}
A billionaire moves into a town of 10,000 people. Predict the change in
mean and in median wealth.
\end{exercise}

\begin{exercise}[$\star\star\star$]\label{exo:b1:describing-data:build}
Department A admits 80 of 100 male and 18 of 20 female applicants;
department B admits 10 of 50 male and 30 of 130 female applicants. Compute
the admission rates by department and overall, and explain.
\end{exercise}

\begin{exercise}[$\star\star\star$]\label{exo:b1:describing-data:minimise}
Show that $\bar x$ minimises $\sum_i (x_i - c)^2$ over $c$, and that a median minimises
$\sum_i |x_i - c|$.
\end{exercise}

\begin{exercise}[$\star\star\star$]\label{exo:b1:describing-data:hidden}
A study reports a correlation of $0.1$ between two variables. What
observation would show that a strong relation lies behind it?
\end{exercise}

\section{Problem: Brenland's Household Survey}\label{sec:b1:describing-data:problem}

\begin{problem}[{Weekend problem --- a survey, described}]\label{pb:b1:describing-data:survey}
Brenland's household survey of 2025 (invented) has 2,000 households, with
disposable and gross income, the head's age, schooling and hours, and the
household's size. Each answer comes from a script.

\medskip\noindent\textbf{Part I --- The distribution.}
\begin{enumerate}
  \item Give the counts of \omterm{def:b1:national-accounts-in-depth:gni}{disposable income} in bins of 10,000 dalens.
  \item Give the counts of $\log_{10}$ of \omterm{def:b1:national-accounts-in-depth:gni}{disposable income} in bins of $0.2$ from $3.8$ to $5$.
  \item Which \omterm{def:b1:describing-data:histogram}{histogram} is closer to symmetric?
  \item Give the mean and the median.
  \item What share of households is below the mean?
  \item Why is the mean above the median?
\end{enumerate}

\medskip\noindent\textbf{Part II --- Moments and quantiles.}
\begin{enumerate}[resume]
  \item Give the standard deviation and the \omterm{def:b1:describing-data:skewness}{skewness}.
  \item Give the \omterm{def:b1:describing-data:skewness}{skewness} of the logarithm.
  \item Give the 10th and 90th \omterm{def:b1:describing-data:quantile}{percentiles} and their ratio.
  \item What share of all \omterm{def:b1:national-accounts-in-depth:gni}{disposable income} goes to the top tenth of households?
  \item Why may a survey understate that share?
\end{enumerate}

\medskip\noindent\textbf{Part III --- Income and schooling.}
\begin{enumerate}[resume]
  \item Give mean gross income for heads with under 12, 12, 13--15 and 16 or
    more years of schooling.
  \item Give the correlation of schooling with gross income, and with its
    logarithm.
  \item Why might the correlation understate how much schooling and income
    are related?
  \item Does the table show that schooling raises income?
\end{enumerate}

\medskip\noindent\textbf{Part IV --- The paradox.} Split households by the
head's age (under 50, 50 or over) and by whether the head works.
\begin{enumerate}[resume]
  \item Give mean \omterm{def:b1:national-accounts-in-depth:gni}{disposable income} of each age group, overall.
  \item Give the four means within working and non-working households.
  \item Give the share working in each age group.
  \item Rebuild each overall mean as a weighted average of the within-group means.
  \item Compute the older households' mean if they worked as often as the younger.
  \item Which comparison answers ``do older households have more income
    when they work?''
  \item Which answers ``do older households have more income?''
  \item Name the trap.
  \item What would a reader of the overall means alone conclude?
  \item State the result.
\end{enumerate}
\end{problem}
